%MB9T5
\setcounter{paragraph}{0}
\setcounter{ex}{0}
\PhanI




\begin{ex}%[3P5Y3-1][Loai3]
Trong thí nghiệm Young về giao thoa ánh sáng, hai khe được chiếu bằng ánh sáng đơn sắc. Nếu tại điểm M trên màn quan sát là vân sáng thì hiệu đường đi của ánh sáng từ hai khe ${S_1},\ {S_2} $ đến M bằng
\choice
{\True số nguyên lần bước sóng}
{số nguyên lần nửa bước sóng}
{bán nguyên lần bước sóng}
{nửa bước sóng}
\loigiai{Tại M là vân sáng: $\Delta d=k\lambda $ với k là số nguyên.
}
\end{ex} \haidong{20}
\begin{ex}%[3P5Y3-1][Loai3]
Trong thí nghiệm Young về giao thoa ánh sáng, hai khe được chiếu bằng ánh sáng đơn sắc. Nếu tại điểm M trên màn quan sát là vân tối thì hiệu đường đi của ánh sáng từ hai khe ${S_1},\ {S_2} $ đến M bằng
\choice
{số nguyên lần bước sóng}
{số nguyên lần nửa bước sóng}
{\True bán nguyên lần bước sóng}
{nửa bước sóng}
\loigiai{Tại M là vân tối: $\Delta d=(k+0,5)\lambda $ với k là số nguyên.
}
\end{ex} \haidong{20}
\loai{Thay ánh sáng khác}
\begin{ex}%[3P5Y3-1][Loai4]
Trong thí nghiệm Young về giao thoa ánh sáng, nếu thay ánh sáng đơn sắc màu lam bằng ánh sáng đơn sắc màu vàng và giữ nguyên các điều kiện khác thì trên màn quan sát
\choice
{\True khoảng vân tăng lên}
{khoảng vân giảm xuống}
{vị trí vân trung tâm thay đổi}
{khoảng vân không thay đổi}
\loigiai{
Ta có: $i=\dfrac{\lambda D}{a} $ mà ${{\lambda}_{l}}<{{\lambda}_{v}}\Rightarrow {i_l}<{i_v} \Rightarrow$ khoảng vân tăng lên.
}
\end{ex} \haidong{20}
\begin{ex}%[3P5Y3-1][Loai7]
Hiện tượng nào sau được ứng dụng để đo bước sóng ánh sáng?
\choice
{\True Hiện tượng giao thoa}
{Hiện tượng quang điện}
{Hiện tượng tán sắc}
{Hiện tượng quang - phát quang}
\loigiai{
Ứng dụng của hiện tượng giao thoa là để đo bước sóng ánh sáng.
}
\end{ex} \haidong{20}
\begin{ex}%[3P5Y3-1][Loai6]
Cho các loại ánh sáng đơn sắc sau: Ánh sáng lục (I); Ánh sáng đỏ (II); Ánh sáng vàng (III); Ánh sáng tím (IV). Loại ánh sáng nào trên hình ảnh giao thoa có khoảng vân lần lượt lớn nhất và nhỏ nhất?
\choice
{I; IV}
{II; III}
{III; IV}
{\True II; IV}
\loigiai{
Ta có: $i=\dfrac{\lambda D}{a} $ mà ${{\lambda}_{II}}>{{\lambda}_{III}}>{{\lambda}_{I}}>{{\lambda}_{IV}}\Rightarrow {{i}_{II}}>{{i}_{III}}>{i_I}>{{i}_{IV}}$.
}
\end{ex} \haidong{20}










\loai{Tính khoảng vân}
\begin{ex}%[3P5-3.2B][Loai1]
\nguon{QG - 2018} Trong thí nghiệm Young về giao thoa ánh sáng, khoảng cách giữa hai khe là 0,5 mm, khoảng cách từ mặt phẳng chứa hai khe đến màn quan sát là 2 m. Chiếu sáng các khe bằng bức xạ có bước sóng 500 nm. Trên màn, khoảng cách giữa hai vân sáng liên tiếp là
\choice
{0,5 mm}
{1 mm}
{4 mm}
{\True 2 mm}
\loigiai{
\begin{itemize}
\item Khoảng cách giữa hai vân sáng liên tiếp: $i=\dfrac{\lambda D}{a}=2\ mm$.
\end{itemize}
}
\end{ex} \haidong{20}

\loai{Từ khoảng vân tìm các đại lượng}
\begin{ex}%[3P5B3-2][Loai2]
\nguon{QG - 2018} Trong thí nghiệm Young về giao thoa ánh sáng, nguồn sáng phát ra ánh sáng đơn sắc có bước sóng 600 nm. Khoảng cách từ mặt phẳng chứa hai khe đến màn quan sát là 2 m. Trên màn, khoảng vân đo được là 1,5 mm. Khoảng cách giữa hai khe bằng
\choice
{0,4 mm}
{0,9 mm}
{0,45 mm}
{\True 0,8 mm}
\loigiai{
\begin{itemize}
\item Khoảng cách giữa hai khe: $a=\dfrac{\lambda D}{i}= 0,8\ mm$.
\end{itemize}
}
\end{ex} \haidong{20}

\loai{Xác định vị trí vân}
\begin{ex}%[3P5B3-2][Loai3]
Thực hiện thí nghiệm Young về giao thoa với ánh sáng đơn sắc có bước sóng 0,4 $\mu m $, khoảng cách giữa hai khe là 0,5 mm, khoảng cách từ mặt phẳng chứa hai khe đến màn là 1 m. Trên màn quan sát, vân sáng bậc 4 cách vân sáng trung tâm là
\choice
{\True 3,2 mm}
{4,8 mm}
{1,6 mm}
{2,4 mm}
\loigiai{Vân sáng bậc 4 cách vân sáng trung tâm: $x=ki=4i=4\dfrac{\lambda D}{a}=4\dfrac{0,4. 1}{0,5}=3,2\ mm $}%<MyLT1>
\end{ex} \haidong{20}

\loai{Từ vị trí vân tìm bước sóng}
\begin{ex}%[3P5B3-2][Loai4]
Trong thí nghiệm giao thoa ánh sáng với hai khe Young cách nhau 0,5 mm, màn quan sát đặt song song với mặt phẳng chứa hai khe và cách hai khe một đoạn 1 m. Tại vị trí M trên màn, cách vân sáng trung tâm một đoạn 4,4 mm là vân tối thứ 6. Bước sóng $\lambda $ của ánh sáng đơn sắc trong thí nghiệm là
\choice
{\True 0,4 $\mu m $}
{0,6 $\mu m $}
{0,5 $\mu m $}
{0,44 $\mu m $}
\loigiai{
\begin{itemize}
\item Ta có: $i=\dfrac{OM}{k}=\dfrac{4,4}{5,5}=0,8\ mm\Rightarrow \lambda =\dfrac{ia}{D}=\dfrac{0,8. 0,5}{1}=0,4\ \mu m $ (vân tối nên k bán nguyên).
\end{itemize}
}%<MyLT1>
\end{ex} \haidong{20}
\loai{Xác định M thuộc vân sáng hay vân tối}
\begin{ex}%[3P5B3-2][Loai5]
Trong thí nghiệm Young về giao thoa ánh sáng, hai khe Young cách nhau 3 mm được chiếu bằng ánh sáng đơn sắc có bước sóng 0,60 $\mu m $. Các vân giao thoa được hứng trên màn đặt cách hai khe 2 m. Tại điểm N cách vân trung tâm 1,8 mm là
\choice
{vân sáng bậc 4}
{vân tối thứ 4}
{\True vân tối thứ 5}
{vân sáng bậc 5}
\loigiai{
\begin{itemize}
\item Khoảng vân $i=\dfrac{\lambda D}{a}=\dfrac{0,6. 2}{3}=0,4\ mm $ $\Rightarrow k=\dfrac{ON}{i}=\dfrac{1,8}{0,4}=4,5 $.
\end{itemize}
}
%<MyLT2>
\end{ex} \haidong{20}

 
\loai{Tính khoảng cách giữa n vân liên tiếp}
\begin{ex}%[3P5K3-2][Loai6]
Ánh sáng chiếu vào hai khe trong thí nghiệm Young là ánh sáng đơn sắc có bước sóng $\lambda$. Tại một điểm M nằm trong vùng giao thoa trên màn cách vân trung tâm là 2,16 mm có hiệu đường đi của ánh sáng từ hai khe đến đó bằng $1,62\ \mu m$. Nếu bước sóng $\lambda =0,6\ \mu m$ thì khoảng cách giữa 5 vân sáng kế tiếp bằng
\choice
{1,6 mm}
{\True 3,2 mm}
{4 mm}
{2 mm}
\loigiai{\begin{itemize}
\item Tại điểm M, ta có:
\begin{align*}
d_2-d_1=\dfrac{ax}{D}\Rightarrow \dfrac{a}{D}=\dfrac{d_2-d_1}{x}=\dfrac{1,62.10^{-6}}{2,16.10^{-3}}=0,75.10^{-3}\Rightarrow \dfrac{D}{a}=\dfrac{4000}{3}.
\end{align*}
\item Khoảng cách giữa 5 vân sáng kế tiếp bằng:
\begin{align*}
\Delta x=\left({5-1}\right)i=4.\dfrac{\lambda D}{a}=4.0,6.10^{-6}.\dfrac{4000}{3}=3,2\ mm.
\end{align*}
\end{itemize}}
%<MyLT>
\end{ex} \haidong{20}
\loai{Cho khoảng cách giữa n vân liên tiếp. Tính các đại lượng}
\begin{ex} %[3P5K3-2][Loai7]
\nguon{QG - 2020} Trong thí nghiệm Young về giao thoa ánh sáng đơn sắc, khoảng cách giữa 5 sáng liên tiếp trên màn quan sát là 3,0 mm. Khoảng vân trên màn là
\choice
{1,2 mm}
{0,6 mm} 
{\True 0,75 mm}
{1,5 mm}
\loigiai{
\begin{itemize}
\item Khoảng vân: $i=\dfrac{3}{4}=0,75\ mm$.
\end{itemize}
}
\end{ex} \haidong{20}



\loai{Xác định khoảng cách giữa hai vân sáng}
\begin{ex}%[3P5K3-2][Loai8]
Trong thí nghiệm Young về giao thoa ánh sáng khoảng cách giữa hai khe là 2 mm, khoảng cách từ hai khe đến màn là 1 m, ánh sáng đơn sắc có bước sóng 0,5 $\mu m $. Khoảng cách giữa vân sáng bậc 3 và vân sáng bậc 5 ở hai bên so với vân sáng trung tâm là
\choice
{0,50 mm}
{0,75 mm}
{1,25 mm}
{\True 2 mm}
\loigiai{\begin{itemize}
\item Khoảng cách giữa vân sáng bậc 3 và vân sáng bậc 5 ở 2 bên vân sáng trung tâm là: $(3+5)i=8i $.
\item Ta lại có $i=\dfrac{\lambda D}{a}=\dfrac{0,5. 1}{2}=0,25\ mm\Rightarrow \Delta x=8i=8. 0,25=2\ mm $.
\end{itemize}
}
%<MyLT2>
\end{ex} \haidong{20}


\loai{Khoảng cách giữa hai vân tối hoặc giữa vân tối với vân sáng}
\begin{ex}%[3P5K3-2][Loai9]
Trong thí nghiệm Young về giao thoa ánh sáng, khoảng cách giữa hai khe là $a=1\ mm$, khoảng cách từ hai khe sáng đến màn $D=3\ m$. Bước sóng ánh sáng dùng trong thí nghiệm $\lambda =0,5\ \mu m$. Khoảng cách từ vân tối thứ 2 đến vân tối thứ 5 cùng bên là
\choice
{12\ mm}
{0,75\ mm}
{\True 4,5\ mm}
{625\ mm}
\loigiai{
\begin{itemize}
\item Khoảng cách giữa vân tối thứ 2 đến vân tối thứ 5 cùng phía với nhau so với vân sáng trung tâm bằng: $4,5i-1,5i=3i=4,5\ mm $.
\end{itemize}
}%<MyLT1>
\end{ex} \haidong{20}

\loai{Từ khoảng cách giữa hai vân tìm các đại lượng}
\begin{ex}%[3P5K3-2][Loai10]
Trong thí nghiệm giao thoa ánh sáng với khe Young, khoảng cách giữa hai khe là 1\ mm, khoảng cách từ hai khe đến màn là 1,5 m. Trên màn, người ta đo khoảng cách từ vân sáng bậc 2 đến vân sáng bậc 7 cùng phía so với vân trung tâm là 4,5\ mm. Bước sóng dùng trong thí nghiệm là
\choice
{$\lambda = 0,4\ \mu m$}
{$\lambda = 0,5\ \mu m$}
{\True $\lambda = 0,6\ \mu m$}
{$\lambda = 0,45\ \mu m$}
\loigiai{
\begin{itemize}
\item Ta có: $x_7-x_2=7\dfrac{\lambda D}{a}-2\dfrac{\lambda D}{a}=5\dfrac{\lambda D}{a}\Rightarrow \lambda =\dfrac{(x_7-x_2)a}{5D}=\dfrac{{4,5.10}^{-3}{.10}^{-3}}{5.1,5}=0,6.10^{-6}\ m $.
\item 
\end{itemize}
}
%<MyLT2>
\end{ex} \haidong{20}

 
\loai{Hiệu đường đi}
\begin{ex}%[3P5B3-2][Loai11]
Trong thí nghiệm Young về giao thoa ánh sáng, nguồn sáng gồm các bức xạ có bước sóng lần lượt là: ${{\lambda}_{1}}=750\ nm $, ${{\lambda}_{2}}=675\ nm $ và ${{\lambda}_{3}}=600\ nm $. Tại điểm M trong vùng giao thoa trên màn mà hiệu khoảng cách đến hai khe bằng 1,5 $\mu m $ có vân sáng của bức xạ
\choice
{${{\lambda}_{2}} $ và ${{\lambda}_{3}} $}
{${{\lambda}_{3}} $}
{\True ${{\lambda}_{1}} $}
{${{\lambda}_{2}} $}
\loigiai{
\begin{itemize}
\item Ta có: ${k_1}=\dfrac{{d_2}-{d_1}}{\lambda}=\dfrac{1. 5}{0,75}=2 $ nguyên vậy có vân sáng của bức xạ ${{\lambda}_{1}}$.
\end{itemize}
}
%<MyLT2>
\end{ex} \haidong{20}
\begin{ex}%[3P5B3-2][Loai11]
Trong thí nghiệm Young về giao thoa ánh sáng, hai khe được chiếu bằng ánh sáng đơn sắc có bước sóng $\lambda $. Nếu tại điểm M trên màn quan sát có vân tối thì hiệu đường đi của ánh sáng từ hai khe đến điểm M có độ lớn nhỏ nhất bằng
\choice
{$\dfrac{\lambda}{4} $}
{$\lambda $}
{\True $\dfrac{\lambda}{2} $}
{2 $\lambda $}
\loigiai{
\begin{itemize}
\item Tại M là vân tối nên: $\Delta d=\left(k+0,5 \right)\lambda \Rightarrow \Delta {{d}_{min}}=0,5\lambda $ khi $k = 1 $.
\end{itemize}
}
%<MyLT>
\end{ex} \haidong{20}
\loai{Cho khoảng bước sóng và hai đoạn. Tìm bước sóng}
\begin{ex} %[3P5K3-2][Loai14]
\nguon{QG - 2019} Tiến hành thí nghiệm Young về giao thoa ánh sáng, nguồn phát ra ánh sáng đơn sắc có bước sóng $\lambda $ ($380\ nm<\lambda <760\ nm$). Khoảng cách giữa hai khe là 1 mm, khoảng cách từ mặt phẳng chứa hai khe đến màn quan sát là 1 m. Trên màn, hai điểm A và B là vị trí hai vân sáng đối xứng nhau qua vân trung tâm, C cũng là vị trí một vân sáng. Biết A, B, C cùng nằm trên một đường thẳng vuông góc với các vân giao thoa, $AB=6,4$ mm và $BC=4$ mm. Giá trị của $\lambda $ bằng
\choice
{700 nm}
{500 nm}
{600 nm}
{\True 400 nm}
\loigiai{
\begin{itemize}
\item Ta có $OA=OB=3,2$ mm. 
\item Tại A, B, C đều là các vân sáng nên ta có: $OC=BC-OB=0,8\ mm=8.10^5\ nm=k\dfrac{\lambda D}{a}=1000.k\lambda $. 
\item Với k nguyên thì chỉ có $\lambda =400\ nm$ thỏa mãn.
 
\end{itemize}
}%<MyLT1>
\end{ex} \haidong{20}
\setcounter{paragraph}{0}
\setcounter{ex}{0}
\PhanII
\begin{ex}
Hiện tượng giao thoa ánh sáng được quan sát trong thí nghiệm Young với hai khe hẹp là minh chứng rõ ràng cho tính chất sóng của ánh sáng. Trong thí nghiệm này, các yếu tố như bước sóng ánh sáng, khoảng cách giữa hai khe, và khoảng cách từ khe đến màn đều ảnh hưởng đến sự tạo thành và đặc điểm của các vân giao thoa. 
\choiceTF[t] 
{\True Khi thay đổi ánh sáng trong thí nghiệm thì khoảng vân cũng thay đổi}
{\True Khi tăng khoảng cách từ màn đến hai khe thì khoảng vân tăng}
{Khi giảm khoảng cách giữa hai khe thì khoảng vân giảm}
{Khi tăng bước sóng của ánh sáng thì khoảng vân giảm}
\loigiai{
\begin{itemchoice}
\itemch
\itemch
\itemch
\itemch
\end{itemchoice}
}
\end{ex} \haidong{20}
\begin{ex}
Trong thí nghiệm giao thoa ánh sáng với hai khe của Young, các vân sáng và vân tối xuất hiện xen kẽ trên màn do sự giao thoa của hai sóng ánh sáng kết hợp.
\choiceTF[t]
{Khoảng vân là khoảng cách giữa hai vân sáng}
{\True Khoảng vân là khoảng cách giữa hai vân tối liên tiếp}
{\True Khoảng vân bằng hai lần khoảng cách giữa một vân sáng và một vân tối liền kề}
{Khoảng cách giữa 3 vân sáng liên tiếp là 3 khoảng vân}
\loigiai{
\begin{itemchoice}
\itemch
\itemch
\itemch
\itemch
\end{itemchoice}
}
\end{ex} \haidong{20}
\begin{ex}
Trong thí nghiệm giao thoa Young về giao thoa ánh sáng có bước sóng $\lambda=0,6\ \mu m$, khoảng cách từ màn tới hai khe là $2 \ m$, khoảng cách giữa hai khe là $a=1 \ mm$.
\choiceTF[t]
{\True Tại vị trí trung tâm của trường giao thoa là vân sáng bậc 0}
{Khoảng vân giao thoa là $i=1,2\ cm$}
{Vị trí cách vân trung tâm $x=2,4\ mm$ là vân sáng bậc 3}
{\True Vị trí cách vân trung tâm $x=3\ mm$ là vân tối bậc 3}
\loigiai{
\begin{itemchoice}
\itemch
\itemch
\itemch
\itemch
\end{itemchoice}
}
\end{ex} \haidong{20}
\begin{ex}
Trong thí nghiệm giao thoa Young về giao thoa ánh sáng có bước sóng $\lambda$, khoảng cách từ màn tới hai khe là $1,5 \ m$, khoảng cách giữa hai khe là $a=1 \ mm$. Biết khoảng cách từ vân trung tâm tới vân sáng bậc 5 là $4,5 \ mm$.
\choiceTF[t]
{\True Khoảng vân là $i=0,9\ mm$}
{Bước sóng giao thoa là $\lambda=6000\ nm$}
{Điểm cách vân trung tâm $3,5 \ mm$ là vân sáng bậc 4}
{Giá sử thay đổi nguồn sáng phát ra bước sóng $\lambda=0,4\ \mu m$ thì khoảng vân tăng 1,5 lần}
\loigiai{
\begin{itemchoice}
\itemch
\itemch
\itemch
\itemch
\end{itemchoice}
}
\end{ex} \haidong{20}
\setcounter{paragraph}{0}
\setcounter{ex}{0}
\PhanIII
\begin{ex}
Trong thí nghiệm Young về giao thoa ánh sáng, khoảng cách giữa hai khe là $1 \ mm$, khoảng cách từ mặt phẳng chứa hai khe đến màn quan sát là $2 \ m$. Nguồn sáng đơn sắc có bước sóng $0,45\ \mu m$. Khoảng vân giao thoa trên màn là bao nhiêu $mm$?\shortans[oly]{0,9}
\loigiai{
0,9
}
\end{ex} \haidong{20}
\begin{ex}
Trong thí nghiệm Young về giao thoa ánh sáng, nguồn $S$ phát bức xạ có bước sóng $600 \ nm$; khoảng cách giữa hai khe là $1,5 \ mm$; khoảng cách từ mặt phẳng chứa hai khe đến màn quan sát là $2,4 \ m$. Dịch chuyển mối hàn của cặp nhiệt điện theo đường song song với mặt phẳng chứa hai khe thì cứ sau một khoảng $\Delta d$ bằng bao nhiêu $mm$ thì kim điện kế lệch nhiều nhất?\shortans[oly]{0,96}
\loigiai{
0,96
}
\end{ex} \haidong{20}
\begin{ex}
Tiến hành thí nghiệm Young về giao thoa ánh sáng với ánh sáng đơn sắc nguồn sáng phát ra ánh sáng đơn sắc có bước sóng $\lambda$. Trên màn quan sát, vân sáng bậc 5 cách vân sáng trung tâm 4,5 mm. Khoảng vân giao thoa trên màn bằng bao nhiêu mm?\shortans[oly]{0,9}
\loigiai{
Vân sáng bậc $5: x=5 i \Rightarrow i=\dfrac{x}{5}=0,9 \ mm$.
Đáp số: 0,9
}
\end{ex} \haidong{20}
\begin{ex}
Trong thí nghiệm Young về giao thoa ánh sáng, khoảng cách giữa hai khe là $2 \ mm$, khoảng cách từ mặt phẳng chứa hai khe đến màn quan sát là $1,2 \ m$. Trên màn, quan sát được 9 vân sáng mà khoảng cách giữa hai vân sáng ngoài cùng là $2,4 \ mm$. Bước sóng của ánh sáng trong thí nghiệm bằng bao nhiêu $\mu m$?\shortans[oly]{0,5}
\loigiai{
\begin{itemize}
\item Khoảng cách 9 vân sáng liên tiếp: $8 i=2,4 \ mm \Rightarrow i=0,3 \ mm$
\item Bước sóng: $\lambda=\dfrac{a \cdot i}{D}=\dfrac{2 \cdot 0,3}{1,2}=0,5\ \mu m$
\end{itemize}
Đáp số: 0,5
}
\end{ex} \haidong{20}
\begin{ex}
Trong thí nghiệm Young về giao thoa ánh sáng, khoảng cách giữa hai khe hẹp là $2 \ mm$, khoảng cách từ mặt phẳng chứa hai khe đến màn là $1,5 \ m$. Người ta đo được khoảng cách giữa vân sáng bậc 2 đến vân sáng bậc 7 nằm khác phía so với vân trung tâm có chiều dài là $3,5 \ mm$. Bước sóng của ánh sáng dùng trong thí nghiệm bằng bao nhiêu $\mu m$ (làm tròn kết quả đến chữ số hàng phần trăm)?
\shortans[oly]{0,52}
\loigiai{
\begin{itemize}
\item Khoảng cách vân sáng bậc 2 và bậc 7 khác phía: $9 i=3,5 \Rightarrow i=\dfrac{7}{18} \ mm$
\item Bước sóng: $\lambda=\dfrac{a \cdot i}{D}=\dfrac{2 \cdot 7}{1,5 \cdot 18}=0,52\ \mu m$
\end{itemize}
Đáp số: 0,52
}
\end{ex} \haidong{20}
\begin{ex}
Trong thí nghiệm Young về giao thoa ánh sáng, hai khe hẹp cách nhau $0,6 \ mm$ và cách màn quan sát $1,2 \ m$. Chiếu sáng các khe bằng ánh sáng đơn sắc có bước sóng $\lambda\ (380 \ nm<\lambda<760 \ nm)$. Trên màn, điểm M cách vân trung tâm $2,3 \ mm$ là vị trí của một vân tối. Giá trị của $\lambda$ bằng bao nhiêu $\mu m$?\shortans[oly]{0,46}
\loigiai{
\begin{itemize}
\item Vị trí vân tối: $x_{t}=(k-0,5) \dfrac{\lambda \cdot D}{a} \Rightarrow \lambda=\dfrac{a \cdot x_{t}}{(k-0,5) \cdot D}=\dfrac{0,6 \cdot 2,3}{(k-0,5) \cdot 1,2}$
\item Bấm table:
\begin{center}
\includegraphics[width=4cm]{images/2024_05_26_49ddb8e3a7584f238c7bg-18}
\end{center}
\end{itemize}
Đáp số: 0,46
}
\end{ex} \haidong{20}


